\documentclass[11pt]{article}
\usepackage[a4paper,margin=2.5cm]{geometry}
\usepackage{amsmath}
\usepackage{amssymb}
\usepackage{enumitem}
\usepackage{parskip}
\setlist[enumerate,2]{label=(\alph*), itemsep=1pt, topsep=2pt}
\setlength{\parindent}{0pt}
\newcommand{\correct}[1]{\textbf{#1}\;$\checkmark$}
\begin{document}
\section*{Social choice}
\begin{enumerate}[itemsep=6pt]
  \item A Condorcet winner is an alternative that ...
  \begin{enumerate}
    \item has the most first choice votes
    \item is ranked last by no one
    \item has the highest Borda score
    \item beats every other alternative in a pairwise majority vote
  \end{enumerate}
  \item A Condorcet cycle is when ...
  \begin{enumerate}
    \item the Borda winner is always unique
    \item pairwise majority preferences cycle even when every voter's ranking is transitive
    \item simple majority rule always picks the Condorcet winner
    \item voters cannot rank candidates
  \end{enumerate}
  \item Borda's method scores each alternative by ...
  \begin{enumerate}
    \item the number of alternatives it is ranked above, summed over all voters
    \item whether it beats every rival head-to-head
    \item a coin toss
    \item the number of first choice votes only
  \end{enumerate}
  \item Arrow's impossibility theorem states that ...
  \begin{enumerate}
    \item simple majority rule is always transitive
    \item Borda's method is the only fair voting rule
    \item no social welfare function satisfies unrestricted domain, Pareto efficiency, independence of irrelevant alternatives and non-dictatorship together
    \item every preference profile has a Condorcet winner
  \end{enumerate}
  \item Simple majority rule ranks \(x\) above \(y\) when ...
  \begin{enumerate}
    \item a strict majority of voters prefer \(x\) to \(y\)
    \item \(x\) has a higher Borda score than \(y\)
    \item any voter prefers \(x\) to \(y\)
    \item \(x\) has more first choice votes than \(y\)
  \end{enumerate}
  \item A social welfare function maps ...
  \begin{enumerate}
    \item every preference profile to a collective ranking of the alternatives
    \item candidates to vote totals only
    \item payoffs to strategies
    \item each voter to a single candidate
  \end{enumerate}
  \item A preference profile is ...
  \begin{enumerate}
    \item the list of every voter's ranking of the alternatives
    \item the winning candidate
    \item a single voter's top choice
    \item a payoff matrix
  \end{enumerate}
  \item Each voter's preference is assumed to be ...
  \begin{enumerate}
    \item a transitive strict ranking of the alternatives
    \item a probability distribution
    \item a single favourite only
    \item possibly cyclic
  \end{enumerate}
  \item Independence of irrelevant alternatives requires that the collective ranking of \(x\) and \(y\) ...
  \begin{enumerate}
    \item always matches the Borda ranking
    \item is decided by a single voter
    \item ignores how voters rank \(x\) and \(y\)
    \item depends only on how voters rank \(x\) against \(y\), not on any third alternative
  \end{enumerate}
  \item When a Condorcet winner exists, it is ...
  \begin{enumerate}
    \item one of several
    \item unique
    \item never the Borda winner
    \item always the Borda winner
  \end{enumerate}
  \item Borda's method always ...
  \begin{enumerate}
    \item produces a complete ranking and never cycles
    \item fails when preferences differ
    \item selects the Condorcet winner
    \item ignores lower preferences
  \end{enumerate}
  \item The Borda winner ...
  \begin{enumerate}
    \item never exists
    \item need not be the Condorcet winner, even when one exists
    \item is always the simple majority winner
    \item is always the Condorcet winner
  \end{enumerate}
  \item A Condorcet cycle has the form ...
  \begin{enumerate}
    \item \(A \succ B \succ C\) with no return
    \item everyone ranking \(A\) first
    \item \(A \succ B \succ C \succ A\) under pairwise majority
    \item a tie between all candidates
  \end{enumerate}
  \item Condorcet's method can fail to select a winner because ...
  \begin{enumerate}
    \item voters cannot rank candidates
    \item it always produces a tie
    \item it ignores first choice votes
    \item pairwise majorities may cycle, so no candidate beats all others
  \end{enumerate}
  \item Different voting methods applied to the same profile ...
  \begin{enumerate}
    \item can select different winners
    \item are forbidden
    \item always tie
    \item always agree
  \end{enumerate}
  \item A voter's preference function is a strict linear order, meaning it is asymmetric, transitive and ...
  \begin{enumerate}
    \item reflexive
    \item random
    \item cyclic
    \item complete (every pair of distinct alternatives is ranked)
  \end{enumerate}
  \item Arrow's impossibility theorem applies to sets of alternatives of size ...
  \begin{enumerate}
    \item exactly 2
    \item at least 10
    \item at least 3
    \item any size, including 1
  \end{enumerate}
  \item Both Borda's method and Condorcet's method respond to Arrow's theorem by giving up which property?
  \begin{enumerate}
    \item independence of irrelevant alternatives
    \item unrestricted domain
    \item non-dictatorship
    \item Pareto efficiency
  \end{enumerate}
  \item With three alternatives, Borda's method awards a voter's top choice ...
  \begin{enumerate}
    \item 2 points
    \item 3 points
    \item 1 point
    \item 0 points
  \end{enumerate}
  \item A voting rule is open to strategic manipulation when a voter can gain by ...
  \begin{enumerate}
    \item adding a new candidate
    \item abstaining from the vote
    \item ranking alternatives honestly
    \item misreporting their true preferences
  \end{enumerate}
\end{enumerate}
\end{document}
